\ifdefined\standalone
\documentclass[12pt]{article}
\usepackage{graphicx}
\usepackage{float}
\usepackage[a4paper, margin=1in]{geometry}
\usepackage[utf8]{inputenc}
\usepackage{textcomp}
\usepackage{amsmath}
\usepackage{booktabs}
\usepackage{tikz}
\usepackage{enumitem}
\usetikzlibrary{decorations.pathmorphing,patterns}
\usepackage[hidelinks]{hyperref}
\usepackage[toc,page]{appendix}

\begin{document}
\fi

\section{Gas species advection}

A one-dimensional porous layer of thickness $ L = 10~\mathrm{cm} $ is considered, with a uniform porosity $ \overline{\psi} = 0.6 $. Gas transport occurs solely by advection, with no diffusion and no volumetric source terms. The gaseous mixture consists of oxygen ($ \mathrm{O_2} $) and nitrogen ($ \mathrm{N_2} $) and has a density $ \rho_g = 1.3~\mathrm{kg/m^3} $. The external mass flux is fixed to $ \dot{m}'' = 10~\mathrm{g.m^{-2}.s^{-1}} $. The entering gas mass fraction is $ Y_{\mathrm{O_2}}^{\infty} = 0.22 $ and $ Y_{\mathrm{N_2}}^{\infty} = 0.78 $. Initially, the porous medium contains only nitrogen, such that $ Y_{0,\mathrm{O_2}} = 0 $ and $ Y_{0,\mathrm{N_2}} = 1 $ throughout the domain. 

Under these assumptions, the transport equation for each gaseous species $ j $ reduces to the advection equation:
\begin{equation}
\overline{\psi} \rho_g \frac{\partial Y_j}{\partial t} + \dot{m}'' \cdot \nabla Y_j = 0
\end{equation}

The top boundary ($ z = 0 $) is impermeable to diffusion:
\begin{equation}
\left. \partial_z Y_j \right|_{z=0} = 0.
\end{equation}

At the bottom boundary ($ z = L $), a Dirichlet boundary condition is imposed:
\begin{equation}
Y_j(L,t) = Y_j^{\infty}.
\end{equation}

This equation admits an analytical solution corresponding to the translation of a step profile at a constant velocity $v = {\dot{m}''}/{\overline{\psi}\rho_g}$. The analytical solution for the species mass fraction can therefore be written as
\begin{equation}
Y_j(z,t) =
\begin{cases}
Y_j^{\infty}, & L-z < v t \\
Y_{0,j}, & L-z \ge v t
\end{cases}
\end{equation}

The numerical simulations are performed using a uniform spatial discretization with $ n_z = 500 $ cells, corresponding to a grid spacing of $ \Delta z = 0.2~\mathrm{mm} $. A fixed time step of $ \Delta t = 0.01~\mathrm{s} $ is used.

Figures~\ref{fig:N2_advection} and~\ref{fig:O2_advection} present the evolution of the $ \mathrm{O_2} $ and $ \mathrm{N_2} $ mass fraction profiles. Numerical results are compared with the analytical solution at several times.

\begin{figure}[H]
\centering
\begin{minipage}{0.49\textwidth}
\includegraphics[width=\textwidth]{N2_advection.png}
\end{minipage}
\hfill
\begin{minipage}{0.49\textwidth}
\includegraphics[width=\textwidth]{O2_advection.png}
\end{minipage}
\caption{Species advection with $ 10~\mathrm{g/m^2/s} $.}
\label{fig:N2_advection}
\label{fig:O2_advection}
\end{figure}

The numerical solution accurately captures the propagation velocity of the front, as predicted by the analytical solution. However, the numerical solution exhibits a progressive smoothing of the front compared to the analytical step profile. Such diffusive behavior is due to the first-order spatial discretization of the advection scheme. This numerical diffusion is expected and decreases with smaller spatial steps. Users should be aware of this behavior.


\ifdefined\standalone
\end{document}
\fi
	
